\documentclass[11pt]{article}
\usepackage[margin=1in]{geometry}
\usepackage{amsmath,amssymb}
\usepackage{hyperref}
\usepackage{natbib}
\usepackage{booktabs}
\usepackage{multirow}

\hypersetup{
    colorlinks=true,
    linkcolor=blue,
    citecolor=blue,
    urlcolor=blue
}

\title{Daily Paper Review: March 13, 2026\\OAKS: Benchmarking Online Adaptation to Continual Knowledge Streams}
\author{Internbot Research Assistant}
\date{March 13, 2026}

\begin{document}
\maketitle

\begin{abstract}
Today's review covers \emph{Can Large Language Models Keep Up? Benchmarking Online Adaptation to Continual Knowledge Streams}~\citep{kim2026oaks}, which introduces OAKS, a benchmark evaluating whether LLMs can track dynamically evolving facts in streaming environments. OAKS comprises two datasets: OAKS-BABI (synthetic, 1.2k questions, 128k tokens) and OAKS-Novel (39 human-curated novels, 870 MCQ, 150k tokens avg). Testing 14 models with multiple adaptation strategies, the results are sobering: even Gemini~3~Pro achieves only 66.3\% on OAKS-B, and all models exhibit a dominant \emph{Volatility} failure mode (32.1\% rate of unnecessary answer changes). RAG, rolling windows, and agentic memory systems (HippoRAG-V2, MemAgent, A-Mem) all fail to robustly solve the problem, revealing that continual knowledge adaptation remains fundamentally unsolved.
\end{abstract}

\section{Paper Summary}

\textbf{Title:} Can Large Language Models Keep Up? Benchmarking Online Adaptation to Continual Knowledge Streams\\
\textbf{Authors:} Jiyeon Kim, Hyunji Lee, Dylan Zhou, Sue Hyun Park, Seunghyun Yoon, Trung Bui, Franck Dernoncourt, Sungmin Cha, Minjoon Seo\\
\textbf{Affiliations:} KAIST, Adobe Research\\
\textbf{Links:}
\href{https://arxiv.org/abs/2603.07392}{arXiv} $\mid$
\href{https://huggingface.co/papers/2603.07392}{HF Daily Papers} $\mid$
\href{https://github.com/kaistAI/OAKS}{GitHub}

\subsection{Core Problem: Static Benchmarks for a Dynamic World}

LLMs operating in real-world contexts encounter knowledge that evolves continuously: facts update, entities change state, new information supersedes old. Existing continual learning benchmarks involve few updates or divergent facts---they do not evaluate fine-grained tracking of the \emph{same fact evolving multiple times} over a streaming context. OAKS fills this gap by measuring interval-level accuracy as facts change across sequential context chunks.

\subsection{Benchmark Design}

\paragraph{OAKS-BABI (OAKS-B).} Synthetic dataset from BABILong, reformulated for dynamic tracking. 1.2k questions over 65 chunks (${\sim}$128k tokens). Average 4.7 answer changes per question. Four question types: tracking, counting, bridge (multi-hop), comparison. Stratified into Sparse/Moderate/Frequent subsets by transition frequency.

\paragraph{OAKS-Novel (OAKS-N).} 39 full-length literary novels with expert human curation (\$17.4k annotation cost). 870 multiple-choice questions (avg 5.5 options), 150.6k tokens per book. 4.7 answer changes per question on average. 55\% retention rate after quality filtering.

\paragraph{Metrics.} Primary: interval-level accuracy (prediction vs.\ ground truth at each time $t$). Detailed: Acquisition Latency (AL, delay to correct prediction after a transition), Distraction Susceptibility (DS, losing a correct answer within a phase), Phase Miss Rate (PM, phases where all predictions are wrong).

\subsection{Methods Evaluated}

\begin{itemize}
    \item \textbf{Context strategies:} Base (concatenate all prior chunks, truncate if needed), RAG (Qwen3-Embedding-0.6B, top-30 chunks), Rolling Window (30 recent chunks), Hybrid RAG+RW.
    \item \textbf{Agentic memory:} HippoRAG-V2, MemAgent, A-Mem.
    \item \textbf{14 models:} Qwen3 (4B--235B), Qwen2.5-7B, GPT-OSS (20B, 120B), Gemma~3 (4B, 27B), Gemini~2.5 (Flash, Pro), Gemini~3~Pro.
\end{itemize}

\subsection{Main Results}

\begin{table}[h]
\centering
\small
\begin{tabular}{lcccc}
\toprule
\textbf{Model} & \textbf{OAKS-B} & \textbf{OAKS-N} & \textbf{B-Sparse} & \textbf{B-Frequent} \\
\midrule
Qwen3-4B & 29.3 & 55.8 & --- & --- \\
Qwen3-30B & 37.8 & --- & --- & --- \\
Qwen3-80B & 43.5 & 65.3 & --- & --- \\
Qwen3-235B & 44.7 & 66.0 & --- & --- \\
GPT-OSS-120B & 33.7 & --- & --- & --- \\
Gemini 2.5 Pro & 62.2 & 75.8 & --- & --- \\
\textbf{Gemini 3 Pro} & \textbf{66.3} & \textbf{75.5} & --- & --- \\
\midrule
\multicolumn{5}{l}{\emph{Open-source avg (by transition freq.)}} \\
--- & --- & --- & 40.4 & 30.2 \\
\multicolumn{5}{l}{\emph{Closed-source avg (by transition freq.)}} \\
--- & --- & --- & 64.2 & 50.3 \\
\bottomrule
\end{tabular}
\end{table}

\paragraph{Thinking modes help substantially.} Qwen3-30B: +7.8pp overall, +16.7pp on tracking questions. Gemini~2.5~Pro: +17.4pp overall, +21.2pp on bridge questions.

\paragraph{Agentic memory systems disappoint.} On Qwen2.5-7B (OAKS-B): Base 24.7\%, RAG 32.4\%, HippoRAG-V2 20.8\%, MemAgent 31.3\%, A-Mem 30.3\%. No agentic system consistently beats simple RAG.

\paragraph{RAG fails on frequent updates.} On the Frequent subset, RAG provides $-$0.8pp vs.\ Base for closed-source models---retrieval noise from outdated chunks actively hurts.

\subsection{Failure Mode Taxonomy}

The paper introduces a behavioral archetype analysis:

\begin{itemize}
    \item When facts \textbf{do change}: \emph{Adaptability} (correct change, 38.1\%), \emph{Maladaptation} (incorrect change, 29.7\%), \emph{Lag} (delayed change, 11.1\%).
    \item When facts \textbf{don't change}: \emph{Volatility} (unnecessary change, \textbf{32.1\%}---the dominant failure), \emph{Stability} (correct, 30.7\%), \emph{Obstinacy} (stuck on wrong answer, 26.1\%).
\end{itemize}

Models split into two behavioral clusters: \emph{change-dominant} (GPT-OSS, Qwen3; 63.2\% change rate $\to$ Volatility/Maladaptation) vs.\ \emph{stay-dominant} (Gemini, Gemma; 55.8\% stay rate $\to$ Obstinacy).

\subsection{Limitations}

\begin{enumerate}
    \item \textbf{Inference-time only:} Evaluates context-based adaptation (RAG, memory), not parametric weight updates (fine-tuning, continual learning algorithms).
    \item \textbf{14 models:} Computational constraints limit breadth; no open-source models above 235B.
    \item \textbf{English only:} Novels and synthetic data are English; multilingual adaptation untested.
    \item \textbf{Literary domain:} OAKS-N uses novels; real-world streams (news, sensor data, conversations) may exhibit different dynamics.
\end{enumerate}

\section{Follow-up Research Directions}

\subsection{SPoT-Style Surgical Adaptation for Continual Knowledge Streams}

\textbf{Gap:} OAKS evaluates only inference-time adaptation (context/retrieval). The benchmark reveals that even 235B models with RAG cap at 66\% accuracy. Parametric adaptation---actually updating model weights as knowledge evolves---is not tested, yet it is the core mechanism studied in continual learning. The challenge is catastrophic forgetting: na\"ive fine-tuning on new facts would overwrite old ones.

\textbf{Approach:} Adapt SPoT's~\citep{lin2026spot} surgical post-training for streaming knowledge. When a new fact arrives that contradicts the model's current prediction: (1)~generate the model's current (wrong) response $y^-$, (2)~minimally edit it to incorporate the new fact, producing $y^+$ that differs only at the fact-bearing tokens, (3)~train with SPoT-BCO's reward-based BCE objective using the pre-update checkpoint as the reference. The Elastic Tether ensures that only the contradicted knowledge is updated while all other knowledge is preserved. The LCS filter ($\gamma{=}0.6$) naturally constrains the scope of each update. Run this as an online loop: for each new context chunk, detect changed facts (via the model's own predictions vs.\ new evidence), surgically update, then proceed.

\textbf{Feasibility:} Moderate (2--3 months). Requires building a fact-change detector (comparing model predictions across intervals) and adapting SPoT's offline pipeline to an online streaming setting. The key research question is whether cumulative surgical updates compound drift over hundreds of intervals or whether the Elastic Tether bounds it---OAKS-B's 65 chunks and 4.7 transitions per question provide a natural testbed.

\subsection{RL-Trained Adaptive Memory Controller}

\textbf{Gap:} The agentic memory systems tested (HippoRAG-V2, MemAgent, A-Mem) use fixed retrieval strategies that fail to adapt to transition frequency. RAG hurts on frequent updates because it retrieves outdated chunks. The missing component is a \emph{learned} controller that decides \emph{what to remember, what to forget, and when to update} based on the knowledge dynamics it observes.

\textbf{Approach:} Train a lightweight memory controller via RL (GRPO) on OAKS-B, where the controller's actions are: (1)~\texttt{store}---add current chunk to memory, (2)~\texttt{update}---replace an existing memory entry, (3)~\texttt{forget}---remove an outdated entry. The reward is OAKS's interval-level accuracy: correct predictions at each interval yield $+1$, incorrect yield $0$. The controller observes the LLM's confidence delta between consecutive intervals as state signal---large confidence shifts suggest a fact transition. Use BandPO's~\citep{li2026bandpo} probability-aware clipping for the controller's policy optimization to maintain exploration of novel memory strategies. Evaluate on OAKS-N (held out during training) to test generalization.

\textbf{Feasibility:} Moderate (3--4 months). The controller is small (can be a separate small LM or MLP); the main challenge is defining the state space and action space precisely enough for RL to learn useful policies. OAKS-B's synthetic structure provides clean reward signals for initial training.

\subsection{Continual Knowledge Streams for Multi-Agent Systems via HACRL}

\textbf{Gap:} OAKS evaluates single models in isolation. Real-world deployment often involves multiple specialized agents with different knowledge bases (e.g., a news agent, a domain expert, a memory agent). HACRL~\citep{zhang2026hacrl} showed that heterogeneous agents can share knowledge via rollout transfer, but this was for static math tasks. Extending to streaming knowledge creates a new challenge: agents have \emph{temporally staggered} knowledge---one agent may have processed a recent update that another has not yet seen.

\textbf{Approach:} Deploy $n$ agents with different context windows or update frequencies on the same OAKS stream. Agent $A$ (large context, slow updates) provides stable baseline knowledge; Agent $B$ (small context, fast updates) tracks recent changes. Use HACRL's cross-agent rollout sharing to transfer knowledge: when $B$ detects a fact transition, it generates rollouts reflecting the new fact and shares them with $A$, using HACPO's capability-aware advantage estimation to weight the update by $B$'s local confidence. This creates a ``temporal knowledge pipeline'' where fast-adapting agents feed updates to slow-adapting ones without destabilizing their broader knowledge.

\textbf{Feasibility:} Hard (4--6 months). Requires implementing multi-agent streaming evaluation on OAKS, defining temporal knowledge transfer protocols, and adapting HACRL's synchronous training to an asynchronous streaming setting. The key research question is whether cross-agent transfer accelerates adaptation without amplifying Volatility errors.

\section{Conclusion}

OAKS reveals that continual knowledge adaptation is a fundamental unsolved challenge for LLMs. Even frontier models (Gemini~3~Pro) achieve only 66.3\% accuracy on dynamically evolving facts, and the dominant failure mode---Volatility (32.1\% unnecessary changes)---shows that models cannot distinguish true knowledge transitions from noise. Agentic memory systems provide no clear advantage over simple RAG, and RAG itself fails when facts change frequently.

For the lab, this paper fills our last uncovered core topic (continual learning for LLM agents) and connects naturally to our review series. Direction~1 (SPoT for streaming knowledge) directly applies our reviewed surgical post-training method to the parametric adaptation gap that OAKS identifies but does not test. Direction~2 (RL-trained memory controller with BandPO) bridges our RL expertise with the memory management challenge. Together, these directions could produce the first methods that actually \emph{solve} OAKS rather than just being benchmarked on it.

\bibliographystyle{plainnat}
\begin{thebibliography}{3}

\bibitem[Kim et~al.(2026)]{kim2026oaks}
J.~Kim, H.~Lee, D.~Zhou, S.H.~Park, S.~Yoon, T.~Bui, F.~Dernoncourt, S.~Cha, and M.~Seo.
\newblock Can large language models keep up? Benchmarking online adaptation to continual knowledge streams.
\newblock \emph{arXiv:2603.07392}, 2026.

\bibitem[Lin \& Han(2026)]{lin2026spot}
W.~Lin and K.~Han.
\newblock Surgical post-training: Cutting errors, keeping knowledge.
\newblock \emph{arXiv:2603.01683}, 2026.

\bibitem[Li et~al.(2026)]{li2026bandpo}
Y.~Li, B.~Wang, Y.~Gao, Y.~Yao, X.~Wang, Z.~Yin, and X.~Qiu.
\newblock BandPO: Bridging trust regions and ratio clipping via probability-aware bounds for LLM reinforcement learning.
\newblock \emph{arXiv:2603.04918}, 2026.

\bibitem[Zhang et~al.(2026)]{zhang2026hacrl}
Z.~Zhang, Z.~Huang, X.~Xia, D.~Wang, F.~Zhuang, S.~Ma, N.~Ding, Y.~Yang, J.~Li, and Y.~Ban.
\newblock Heterogeneous agent collaborative reinforcement learning.
\newblock \emph{arXiv:2603.02604}, 2026.

\end{thebibliography}

\end{document}
